\chapter{六阶对称不定矩阵的延迟主元分解示例}

本附录通过一个 $6\times6$ 对称不定矩阵说明延迟主元在多波前分解中的完整传播过程，包括结构排序、消去树构造、波前装配、延迟变量上传、根波前扩展以及 $2\times2$ 主元消去。

\section{原始矩阵与结构排序}

取

\[
\varepsilon=10^{-12},
\qquad
\delta=10^{-14},
\]

并考虑对称矩阵

\[
A=
\begin{bmatrix}
\varepsilon&\delta&1&0&0&0\\
\delta&4&0&1&0&0\\
1&0&0&3&1&0\\
0&1&3&9&0&1\\
0&0&1&0&5&0.5\\
0&0&0&1&0.5&6
\end{bmatrix}.
\]

矩阵的关键特征是

\[
a_{11}=\varepsilon=10^{-12},
\qquad
a_{13}=1.
\]

变量 $x_1$ 的对角元素很小，但它与 $x_3$ 具有较强耦合。删除变量集合 $\{x_3,x_4\}$ 后，结构图分为左右两个子域，因此可将 $\{x_3,x_4\}$ 作为分隔集，并采用排列

\[
p=[x_1,x_2,x_5,x_6,x_3,x_4].
\]

该排列先处理左子域 $\{x_1,x_2\}$ 和右子域 $\{x_5,x_6\}$，最后处理分隔集 $\{x_3,x_4\}$。排列后的矩阵为

\[
P^TAP=
\begin{bmatrix}
\varepsilon&\delta&0&0&1&0\\
\delta&4&0&0&0&1\\
0&0&5&0.5&1&0\\
0&0&0.5&6&0&1\\
1&0&1&0&0&3\\
0&1&0&1&3&9
\end{bmatrix}.
\]

\section{消去树与波前树}

消去 $x_1$ 时，其尚未消去的邻接变量为 $\{x_2,x_3\}$，符号消元产生填充边 $x_2$--$x_3$；消去 $x_5$ 时，其尚未消去的邻接变量为 $\{x_6,x_3\}$，产生填充边 $x_6$--$x_3$。由此得到依赖关系

\[
x_1\longrightarrow x_2\longrightarrow x_3\longrightarrow x_4,
\qquad
x_5\longrightarrow x_6\longrightarrow x_3\longrightarrow x_4.
\]

相应的多波前树可组织为

\begin{figure}[htbp]
\centering
\begin{tikzpicture}[
  font=\small,
  box/.style={draw,thick,rounded corners=2pt,minimum width=3.1cm,minimum height=0.85cm,align=center},
  arr/.style={-{Stealth[length=2.4mm]},thick}
]
  \node[box,fill=MumpsBlue!13] (root) at (0,3.8) {$F_{34}$\\候选变量 $\{x_3,x_4\}$};
  \node[box,fill=MumpsCyan!11] (f2) at (-2.3,2.25) {$F_2$\\候选变量 $\{x_2\}$};
  \node[box,fill=MumpsCyan!11] (f6) at (2.3,2.25) {$F_6$\\候选变量 $\{x_6\}$};
  \node[box] (f1) at (-2.3,0.7) {$F_1$\\候选变量 $\{x_1\}$};
  \node[box] (f5) at (2.3,0.7) {$F_5$\\候选变量 $\{x_5\}$};
  \draw[arr] (f1)--(f2);
  \draw[arr] (f2)--(root);
  \draw[arr] (f5)--(f6);
  \draw[arr] (f6)--(root);
\end{tikzpicture}
\caption{六阶算例的多波前树。贡献块由叶节点向根节点传递。}
\label{fig:six-by-six-front-tree}
\end{figure}

其中根波前 $F_{34}$ 负责分隔集，左右两个分支分别处理两个局部子域。

\section{左侧子树中的主元延迟}

\subsection{$F_1$：$x_1$ 无法作为稳定主元}

叶节点 $F_1$ 的变量顺序为

\[
[\,\underbrace{x_1}_{\text{完全求和变量}}
\mid
\underbrace{x_2,x_3}_{\text{更新变量}}\,],
\]

其波前矩阵为

\[
F_1=
\begin{bmatrix}
\varepsilon&\delta&1\\
\delta&0&0\\
1&0&0
\end{bmatrix}.
\]

设主元接受阈值为 $u=0.1$。对 $x_1$ 检查

\[
|a_{11}|=10^{-12},
\qquad
u\max(|\delta|,1)=0.1.
\]

由于

\[
10^{-12}<0.1,
\]

$x_1$ 不能在 $F_1$ 中作为稳定的 $1\times1$ 主元。该节点没有其他完全求和变量可与之组成 $2\times2$ 主元，因此 $x_1$ 被延迟，随贡献信息传递给父节点 $F_2$。此时

\[
NASS=1,
\qquad
NPIV=0,
\qquad
NELIM=1.
\]

\subsection{$F_2$：消去 $x_2$，继续延迟 $x_1$}

$F_2$ 接收延迟变量 $x_1$ 后，完全求和候选集合由 $\{x_2\}$ 扩展为 $\{x_1,x_2\}$。按

\[
[\,\underbrace{x_1,x_2}_{\text{完全求和变量}}
\mid
\underbrace{x_3,x_4}_{\text{更新变量}}\,]
\]

排列，有

\[
F_2=
\begin{bmatrix}
\varepsilon&\delta&1&0\\
\delta&4&0&1\\
1&0&0&0\\
0&1&0&0
\end{bmatrix}.
\]

$x_1$ 的 $1\times1$ 主元仍不能通过阈值检验。若尝试使用 $\{x_1,x_2\}$ 形成二阶主元块，则

\[
B_{12}=
\begin{bmatrix}
\varepsilon&\delta\\
\delta&4
\end{bmatrix},
\qquad
\det(B_{12})=4\varepsilon-\delta^2
\approx4\times10^{-12}.
\]

该块仍然过于接近奇异，不能作为稳定的 $2\times2$ 主元。相比之下，$x_2$ 满足

\[
4\ge0.1\max(|\delta|,1),
\]

因此采用

\[
D_2=[4]
\]

作为 $1\times1$ 主元，消去 $x_2$，而 $x_1$ 继续延迟。对剩余变量 $[x_1,x_3,x_4]$ 形成的左侧贡献块为

\[
C_{\mathrm{left}}=
\begin{bmatrix}
\varepsilon-\delta^2/4&1&-\delta/4\\
1&0&0\\
-\delta/4&0&-1/4
\end{bmatrix}
\approx
\begin{bmatrix}
10^{-12}&1&0\\
1&0&0\\
0&0&-0.25
\end{bmatrix}.
\]

此时

\[
NASS=2,
\qquad
NPIV=1,
\qquad
NELIM=1.
\]

$x_1$ 已跨过 $F_1$ 和 $F_2$ 两个波前，但仍未被消去。

\section{右侧子树的正常消元}

右侧叶节点 $F_5$ 按 $[x_5\mid x_6,x_3]$ 排列：

\[
F_5=
\begin{bmatrix}
5&0.5&1\\
0.5&0&0\\
1&0&0
\end{bmatrix}.
\]

$x_5$ 可以直接作为主元，取

\[
D_5=[5],
\]

得到变量 $[x_6,x_3]$ 上的贡献块

\[
C_5=
\begin{bmatrix}
-0.05&-0.1\\
-0.1&-0.2
\end{bmatrix}.
\]

父节点 $F_6$ 按 $[x_6\mid x_3,x_4]$ 装配为

\[
F_6=
\begin{bmatrix}
5.95&-0.1&1\\
-0.1&-0.2&0\\
1&0&0
\end{bmatrix}.
\]

$x_6$ 同样可以直接作为 $1\times1$ 主元：

\[
D_6=[5.95].
\]

消去 $x_6$ 后，右侧子树向根节点发送

\[
C_{\mathrm{right}}=
\begin{bmatrix}
-0.20168067&0.01680672\\
0.01680672&-0.16806723
\end{bmatrix},
\]

其变量索引为 $[x_3,x_4]$。右侧子树没有产生延迟主元。

\section{根波前的动态扩展与二阶主元}

根节点原计划只处理分隔集 $\{x_3,x_4\}$。由于左侧子树传来了延迟变量 $x_1$，根节点接收的数据变为

\begin{itemize}
  \item 左侧贡献块：变量集合 $\{x_1,x_3,x_4\}$；
  \item 右侧贡献块：变量集合 $\{x_3,x_4\}$；
  \item 分隔集的原始矩阵块：
  \[
  \begin{bmatrix}
  0&3\\
  3&9
  \end{bmatrix}.
  \]
\end{itemize}

执行扩展累加后，根波前按 $[x_1,x_3,x_4]$ 排列为

\[
F_{34}=
\begin{bmatrix}
10^{-12}&1&-2.5\times10^{-15}\\
1&-0.20168067&3.01680672\\
-2.5\times10^{-15}&3.01680672&8.58193277
\end{bmatrix}.
\]

因此，符号分析预计的 $2\times2$ 根波前在数值阶段扩展为 $3\times3$ 波前。

在根节点中，$x_1$ 的 $1\times1$ 主元仍然失败：

\[
10^{-12}<0.1\times1.
\]

$x_3$ 的对角元素也不足以通过当前检验：

\[
0.20168067<0.1\times3.01680672.
\]

但是变量 $\{x_1,x_3\}$ 可以组成主元块

\[
D_{13}=
\begin{bmatrix}
10^{-12}&1\\
1&-0.20168067
\end{bmatrix},
\qquad
\det(D_{13})\approx-1.
\]

该 $2\times2$ 块非奇异且数值稳定，因此作为一个块主元被接受。延迟变量 $x_1$ 最终与根节点变量 $x_3$ 一同消去。最后剩余的 $x_4$ 形成

\[
D_4\approx[8.58193277].
\]

这个过程说明，小对角元素或零对角元素并不等价于矩阵奇异。当前节点找不到稳定的 $1\times1$ 主元时，可以先延迟变量，在更大的候选集合中寻找稳定的 $2\times2$ 主元。

\section{最终数值主元顺序与分解结果}

符号排序给出的计划顺序为

\[
[x_1,x_2,x_5,x_6,x_3,x_4],
\]

实际数值主元顺序则为

\[
[x_2,x_5,x_6,\{x_1,x_3\},x_4].
\]

令

\[
q=[x_2,x_5,x_6,x_1,x_3,x_4],
\]

$Q$ 为对应的对称置换矩阵，则

\[
Q^TAQ=LDL^T.
\]

各个 $1\times1$ 主元生成的主要消元乘子为

\begin{align*}
L_{x_1,x_2}&=\delta/4=2.5\times10^{-15},
&L_{x_4,x_2}&=1/4=0.25,\\
L_{x_6,x_5}&=0.5/5=0.1,
&L_{x_3,x_5}&=1/5=0.2,\\
L_{x_3,x_6}&=-0.1/5.95=-0.01680672,
&L_{x_4,x_6}&=1/5.95=0.16806723.
\end{align*}

对于根节点的 $2\times2$ 主元，$x_4$ 与 $\{x_1,x_3\}$ 的耦合行为

\[
[r,s]=[-2.5\times10^{-15},\,3.01680672].
\]

因此需要执行块求解

\[
[L_{x_4,x_1},L_{x_4,x_3}]
=[r,s]D_{13}^{-1},
\]

得到

\[
L_{x_4,x_1}\approx3.01680672,
\qquad
L_{x_4,x_3}\approx-3.0193\times10^{-12}.
\]

按照实际数值主元顺序 $q$，最终因子近似为

\[
\resizebox{0.93\textwidth}{!}{$
L=
\begin{bmatrix}
1&0&0&0&0&0\\
0&1&0&0&0&0\\
0&0.1&1&0&0&0\\
2.5\times10^{-15}&0&0&1&0&0\\
0&0.2&-0.01680672&0&1&0\\
0.25&0&0.16806723&3.01680672&-3.0193\times10^{-12}&1
\end{bmatrix}
$},
\]

\[
\resizebox{0.68\textwidth}{!}{$
D=
\begin{bmatrix}
4&0&0&0&0&0\\
0&5&0&0&0&0\\
0&0&5.95&0&0&0\\
0&0&0&10^{-12}&1&0\\
0&0&0&1&-0.20168067&0\\
0&0&0&0&0&8.58193277
\end{bmatrix}
$}.
\]

由于矩阵元素和因子在此处按有限位小数展示，上式是舍入后的分解结果。该算例的核心过程可以概括为

\[
\boxed{
\begin{gathered}
x_1\text{ 在 }F_1\text{ 延迟}
\longrightarrow
x_1\text{ 在 }F_2\text{ 继续延迟}\\
\longrightarrow
F_{34}\text{ 动态扩展}
\longrightarrow
\{x_1,x_3\}\text{ 形成稳定的 }2\times2\text{ 主元}
\end{gathered}}
\]
